\section{Synthesis and Crystal-Growth Routes}

\subsection{Solid-state and combustion synthesis}

Solid-state reaction begins with oxide, carbonate, or nitrate precursors and relies on diffusion across reacting particles. Carbonate decomposition and local mixing can create transient composition gradients before a silicate phase appears \cite{cultrone2001carbonate,takesue2009thermal,shen2019synthesis}. Combustion routes reduce the diffusion length and can yield fine powders, but the rapid exotherm and gas release change the local oxygen and temperature history \cite{rakov2011near,gupta2021combustion,richhariya2021synthesis}. These methods are appropriate for screening reaction networks; they do not by themselves establish equilibrium.

\subsection{Flux and solution crystal growth}

Flux growth is the closest analogue to a route for isolating a crystal. Ba--Y trisilicate was grown from a molybdate-based flux, demonstrating that a liquid can enable crystal growth at temperatures below a nominal melt point \cite{kolitsch2006bay2si3o10}. Reviews of quaternary oxide flux growth stress that flux composition, solubility, and supersaturation must be measured rather than copied between systems \cite{bugaris2012materials,kumar1974growth,smith1974flux,wanklyn1983supersaturation}. Recent rare-earth silicate chloride growth further illustrates how flux and mixed-anion chemistry determine the accessible structure family \cite{goncalves2025flux}.

\subsection{Solution, gel, and hydrothermal conversion}

Sol--gel and hydrothermal routes homogenize cations before crystallization and can decouple nucleation from high-temperature solid-state diffusion \cite{ranganathan2008sol,parsa2018synthesis,peng2016hydrothermal,ranganathan2008sol}. In glass-ceramic processing, a parent glass first relaxes and then crystallizes during a controlled heat treatment; in-situ studies of barium zinc silicate show that surface and volume crystallization can coexist \cite{bocker2012in,bai2020kinetic,waurischk2020crystal}. Such routes are useful for probing metastable pathways, but the resulting phase may be history-dependent.

\subsection{Thermal profile and processing boundary}

The minimum route record should include precursor identities and ratios, mixing protocol, flux or solvent, crucible material, atmosphere, heating rate, dwell temperature and time, cooling rate, and recovery procedure. Pt, Au, or Pt--Ir containment is a route variable because it can limit contamination while also acting as a nucleating surface or interacting with a melt \cite{thieme2022noble,abdeyazdan2024phase,abdeyazdan2023experimental}. Glass-crystallization studies show that melting conditions alter viscosity and nucleation tendencies, so two nominally identical schedules can produce different phase fractions \cite{jensen2026impact,hu2025crystallization,shaaban2021thermal}. A route comparison is therefore meaningful only when these conditions are reported together.
